\bibliographystyle{unsrt}
\begin{thebibliography}{99}

\bibitem{KW}
A. Kapustin and E. Witten, ``Electric-Magnetic Duality And The Geometric Langlands Program,'' Commun. Numb. Th. Physics {\bf 1} (2007) 1-236,
hep-th/0604151.

\bibitem{frenkel}
E. Frenkel, ``Lectures On The Langlands Program And Conformal Field Theory,'' arXiv:hep-th/0512172.

\bibitem{GW}
S. Gukov and E. Witten, ``Gauge Theory, Ramification, and The Geometric Langlands Program,'' in J. David et. al. (eds.) {\it Current Developments
In Mathematics, 2006} (International Press, Somerville MA, 2008), 35-180, hep-th/0612073.

\bibitem{Wild}
E. Witten,  ``Gauge Theory And Wild Ramification,'' Analysis and Applications {\bf 6} (2008) 429-501.

\bibitem{FW}
E. Frenkel and E. Witten, ``Geometric Endoscopy and Mirror Symmetry,''  Commun. Num.Theor.Phys. {\bf 2} (2008) 113-283,  arXiv:0710.5939.

\bibitem{GaiW}
D. Gaiotto and E. Witten, ``Supersymmetric Boundary Conditons In $\N=4$ Super Yang-Mills Theory,''
arXiv:0804.2902.

\bibitem{GaiW2}
D. Gaiotto and E. Witten, ``$S$-Duality Of Boundary Conditions In $\N=4$ Super Yang-Mills Theory,'' arXiv:0807.3720.

\bibitem{WittenKhov}
E. Witten, ``Fivebranes And Knots,'' Quantum Topol. {\bf 3} (2012) 1-137.

\bibitem{EW1}
E. Witten, ``Mirror Symmetry, Hitchin's Equations, and Langlands Duality,'' arXiv:0802.0999.

\bibitem{EW2}
E. Witten, ``Geometric Langlands From Six Dimensions,'' arXiv:0905.2720.

\bibitem{EW3}
E. Witten, ``Geometric Langlands And The Equations Of Nahm And Bogomolny,''
arXiv:0905.4795.

\bibitem{BeD}
A. Beilinson and V. Drinfeld,``Quantization Of Hitchin's Integrable System and Hecke Eigensheaves,'' preprint (ca. 1995).

\bibitem{Hitchin}
N. Hitchin, ``The Self-Duality Equations On A Riemann Surface,''
Proc. London Math. Soc. (3) {\bf 55} (1987) 59-126.

\bibitem{Freed}
D. Freed, ``Higher Algebraic Structures And Quantization,'' ” Commun. Math. Phys. {\bf 159}
(1994) 343, arXiv:hep-th/9212115.

\bibitem{BD} J. Baez and J. Dolan, “Higher Dimensional Algebra And Topological Quantum Field Theory,”
J. Math. Phys. {\bf  36} (1995), 6073-6105, arXiv:q-alg/9503002.

\bibitem{Kapustin} A. Kapustin, ``Topological Field Theory, Higher Categories, And Their Applications,''
arXiv:1004.2307.

\bibitem{BJSV}
M. Bershadsky, A. Johanssen, V. Sadov, and C. Vafa, ``Topological Reduction Of $4-D$ SYM to $2-D$ Sigma Models,''
Nucl. Phys. {\bf B448} (1995) 166-86, arXiv:hep-th/9501096.

\bibitem{HMS}
J. A. Harvey, G. W. Moore, and A. Strominger, ``Reducing $S$ Duality To $T$ Duality,'' Phys Rev {\bf D52} (1995) 7161-7.

\bibitem{AB} M. F. Atiyah and R. Bott, “The Yang-Mills Equations Over Riemann Surfaces,” Phil. Trans.
R. Soc. London A308 (1982) 523-615.

\bibitem{Corlette}
K. Corlette, ``Flat $G$-Bundles With Canonical Metric,'' J. Diff. Geom. {\bf 28} (1988) 361-82.

\bibitem{HitchinDuke}
N. Hitchin, ``Stable Bundles And Integrable Systems,'' Duke Math. J. {\bf 54} (1987) 91-114.

\bibitem{DonagiMarkman}
R. Donagi and E. Markman, ``Spectral Curves, Algebraically Completely Integrable Hamiltonian Systems,
and Moduli of Bundles,''  arXiv:alg-geom/9507017.

\bibitem{donga}
R. Donagi and D. Gaitsgory, ``The Gerbe Of Higgs Bundles,'' arXiv:math/0005132.

\bibitem{HauselThaddeus}
T. Hausel and M. Thaddeus, ``Mirror Symmetry, Langlands Duality, and the Hitchin System,''  arXiv:math/0205236.

\bibitem{SYZ}
A. Strominger, S.-T. Yau, and E. Zaslow, ``Mirror Symmetry is $T$-Duality,'' Nucl.  Phys.{\bf  B479} :(1996)  243–259.

\bibitem{KP}
L. Katzarkov and T. Pantev, ``Stable $G_2$ Bundles And Algebraically Completely Integrable Systems,'' Compositio Math. {\bf 92} (1994) 43-60.

\bibitem{Hthree}
N. Hitchin, ``Langlands Duality And $G_2$ Spectral Curves,'' arXiv:math/0611524.

\bibitem{rsw}
 T. R. Ramadas, I. M. Singer, and J. Weitsman, “Some Comments On Chern-Simons
Theory,” Commun. Math. Phys. {\bf 126} (1989) 409.

\bibitem{axelrod}
S. Axelrod, ``Geometric Quantization Of Chern-Simons Gauge Theory,'' Ph.D. thesis, Princeton University (1991).

\bibitem{Kapustintwo}
A. Kapustin, ``Wilson-'t Hooft Operators In Four-Dimensional Gauge Theories and $S$-Duality,''
Phys. Rev. D 74, 025005 (2006),
arXiv:hep-th/0501015.

\bibitem{kronheimermon}Kronheimer, P. B. (1986), MSc. thesis
 (Oxford University), unpublished.

\bibitem{polo}Weinberg, E. J. (1982), ``A Continuous Family Of
Magnetic Monopole Solutions,'' Phys. Lett. {\bf B119}, 151.  .

\bibitem{golo}Lee, K.-M., E. J. Weinberg, and P. Yi (1996), ``Massive And
Massless Monopoles With Nonabelian Magnetic Charges, Phys. Rev.
{\bf D54}, 6351-6371.

\bibitem{golo2}Weinberg, E. J., and P. Yi (1998), ``Explicit
Multimonopole Solutions In $SU(N)$ Gauge
Theory,'' Phys. Rev. {\bf D58}, 046001, hep-th/9803164.

\bibitem{latwein}Houghton, C. J., and E. J. Weinberg (2002), ``Multicloud
Solutions With Massless And Massive Monopoles,'' Phys. Rev. {\bf
D66}, 125002, hep-th/0207141.

\bibitem{bais}Bais, F. A., and B. J. Schroers (1998), ``Quantization Of
Monopoles With Nonabelian Magnetic Charge,'' Nucl. Phys. {\bf
B512}, 250, hep-th/9708004.

\bibitem{hanwit}Hanany, A., and E. Witten (1997), ``Type IIB
Superstrings, BPS Monopoles, And Three-Dimensional Gauge
Dynamics,'' Nucl. Phys. {\bf B492}, 152-190.

\bibitem{cherkis}Cherkis, S., and A. Kapustin (2003), ``Periodic Monopoles With
Singularities And $N=2$ Super QCD,'' Commun. Math. Phys. {\bf 234}, 1-35, hep-th/0011081.

\bibitem{donagip} R. Donagi, private communication (ca. 1995).

\bibitem{deligne}
P. Deligne, ``La Determinant de la Cohomologie,'' Contemporary Mathematics,  {\bf 67} (1987) 93-177.

\bibitem{quillen}D. Quillen, ``Determinants Of Cauchy-Riemann Operators Over A Riemann Surface,''  
Funct. Anal. Appl. {\bf 19} (1985) 31-4.

\end{thebibliography}
